\chapter{Supporting Derivations}\label{ap:supportingDerivations}

\section{Feynman Rules}\label{sec:FeynmanRules}

For the massless final-final quark-antiquark antenna functions considered throughout this 
work, in CDR, with $d=4-2\epsilon$, and using the Feynman gauge $\lambda=1$, the required 
Feynman rules go as follows.

An application of these rules is shown in full in Subsection~\ref{sec:matDerivation}.

\subsection{Propagators}

Each internal line in a Feynman diagram represents a propagator: the way a quantum 
field carries an excitation from one spacetime point to another. These propagators 
differ particle-by-particle both in mathematical structure and in the way they are 
graphically represented in the diagrams. A description of propagators encountered 
throughout this work is present in Table~\ref{tab:props} \cite{Ellis:1996mzs,Peskin:1995ev}.

\begin{table}[h!]
\footnotesize
\centering
\caption{\centering Overview of the propagators found throughout the thesis, alongside their mathematical and graphical representations.}
\renewcommand{\arraystretch}{1.3}
  \begin{tabular}{|Sc|Sc|Sc|} 
   \hline
   \textbf{Particle-type} & \textbf{Mathematical Representation} & \textbf{Graphical Representation} \\ 
   \hline
   Quark ($q$) & $\begin{gathered}\frac{i\delta_{ij}\slashed p}{p^2+i0}\\\left(m_q\neq0:\,\frac{i(\slashed p+m_q)}{p^2-m_q^2+i0}\right)\end{gathered}$ & \begin{tikzpicture}
                                                    \begin{feynman}
                                                        \vertex (a) at (0,0);
                                                        \vertex (b) at (2,0);
                                                        \diagram* {
                                                        (a) -- [fermion] (b),
                                                        };
                                                    \end{feynman}
                                                \end{tikzpicture} \\ \hline
   Gluon ($g$) & $\begin{gathered}\frac{i\delta^{ab}}{p^2+i0}\left(-g_{\mu\nu}+(1-\lambda)\frac{p_\mu p_\nu}{p^2}\right)\\\left(\lambda=1:\,\frac{-i\delta^{ab}g_{\mu\nu}}{p^2+i0}\right)\end{gathered}$ & \begin{tikzpicture}
                                                    \begin{feynman}
                                                        \vertex (a) at (0,0);
                                                        \vertex (b) at (2,0);
                                                        \diagram* {
                                                        (a) -- [gluon] (b),
                                                        };
                                                    \end{feynman}
                                                \end{tikzpicture} \\ \hline
   Ghost ($c$) & $\begin{gathered}\frac{i\delta^{ab}}{p^2+i0}\end{gathered}$            & \begin{tikzpicture}
                                                    \begin{feynman}
                                                        \vertex (a) at (0,0);
                                                        \vertex (b) at (2,0);
                                                        \diagram* {
                                                        (a) -- [ghost] (b),
                                                        };
                                                    \end{feynman}
                                                \end{tikzpicture} \\ \hline
  \end{tabular}
   
\label{tab:props}
\end{table}

For an external massless quark, 
\begin{equation}
    \sum_s u_s(p)\bar u_s(p) = \slashed p,\quad \sum_sv_s(p)\bar v_s(p) = \slashed p.
\end{equation}
For external gluons,
\begin{equation}
    \sum_\lambda\varepsilon_\mu^{(\lambda)}(p)\varepsilon_\nu^{(\lambda)\,*}(p)=-g_{\mu\nu} + t_g,\quad t_g = \frac{p_\mu n_\nu+n_\mu p_\nu}{p\cdot n}
\end{equation}
where $n$ is a light-like reference vector and $t_g$ removes the unphysical (timelike and
longitudinal) gluon polarisations. With a single external gluon ($A_3^0$ and the $A_3^1$ set), $t_g$
drops out by the Ward identity and $-g_{\mu\nu}$ alone is used. With two external gluons ($A_4^0$
and $\tilde A_4^0$) it does not drop out, and \textit{AntCalc} keeps it, taking each gluon's
reference vector to be the other gluon's momentum. 

\subsection{Vertices}

These lines meet in vertices, which represent the way in which quantum fields 
interact. The four vertices encountered throughout this work are shown in 
Table~\ref{tab:vertices} \cite{Ellis:1996mzs,Peskin:1995ev}.

\begin{table}[h!]
\footnotesize
\centering
\caption{\centering Overview of the vertices encountered throughout this thesis, alongside their mathematical and graphical representations.}
  \begin{tabular}{|Sc|Sc|Sc|} 
   \hline
   \textbf{Interaction} & \textbf{Mathematical Representation} & \textbf{Graphical Representation} \\ 
   \hline
   $q\bar qg$ & $\displaystyle ig_s\gamma^\mu\left(T^a\right)_{ij}$ & \begin{tikzpicture}[baseline=(current bounding box.center)]
                                                    \begin{feynman}
                                                        \vertex (a) at (0,0);
                                                        \vertex (b) at (0.8,0);
                                                        \vertex (c) at (1.5,0.5);
                                                        \vertex (d) at (1.5,-0.5);
                                                        \diagram* {
                                                        (a) -- [fermion] (b),
                                                        (b) -- [fermion] (c),
                                                        (b) -- [gluon] (d)
                                                        };
                                                    \end{feynman}
                                                \end{tikzpicture} \\ \hline
   $ggg$ & $\begin{gathered}g_sf^{abc}\left(g^{\mu\nu}(p-q)^\rho+g^{\nu\rho}(q-r)^\mu+g^{\rho\mu}(r-p)^\nu\right)\\p+q+r=0\end{gathered}$ & \begin{tikzpicture}[baseline=(current bounding box.center)]
                                                    \begin{feynman}
                                                        \vertex (a) at (0,0);
                                                        \vertex (b) at (0.8,0);
                                                        \vertex (c) at (1.5,0.5);
                                                        \vertex (d) at (1.5,-0.5);
                                                        \diagram* {
                                                        (a) -- [gluon] (b),
                                                        (b) -- [gluon] (c),
                                                        (b) -- [gluon] (d)
                                                        };
                                                    \end{feynman}
                                                \end{tikzpicture} \\ \hline
   $gggg$ & $\begin{aligned} -ig_s^2\Big(f^{abe}&f^{cde}(g^{\mu\rho}g^{\nu\sigma}-g^{\mu\sigma}g^{\nu\rho})\\&+f^{ace}f^{bde}(g^{\mu\nu}g^{\rho\sigma}-g^{\mu\sigma}g^{\nu\rho})\\&+f^{ade}f^{bce}(g^{\mu\nu}g^{\rho\sigma}-g^{\mu\rho}g^{\nu\sigma})\Big)\end{aligned}$            & \begin{tikzpicture}[baseline=(current bounding box.center)]
                                                    \begin{feynman}
                                                        \vertex (a) at (0,0);
                                                        \vertex (b) at (0.5,0.5);
                                                        \vertex (c) at (0.5,-0.5);
                                                        \vertex (d) at (-0.5,-0.5);
                                                        \vertex (e) at (-0.5,0.5);
                                                        \diagram* {
                                                        (a) -- [gluon] (b),
                                                        (a) -- [gluon] (c),
                                                        (a) -- [gluon] (d),
                                                        (a) -- [gluon] (e)
                                                        };
                                                    \end{feynman}
                                                \end{tikzpicture} \\ \hline
   $c\bar cg$ & $\begin{gathered}-g_sf^{abc}p^\mu\\p\text{ is the outgoing ghost momentum}\end{gathered}$            & \begin{tikzpicture}[baseline=(current bounding box.center)]
                                                    \begin{feynman}
                                                        \vertex (a) at (0,0);
                                                        \vertex (b) at (0.8,0);
                                                        \vertex (c) at (1.5,0.5);
                                                        \vertex (d) at (1.5,-0.5);
                                                        \diagram* {
                                                        (a) -- [ghost] (b),
                                                        (b) -- [ghost] (c),
                                                        (b) -- [gluon] (d)
                                                        };
                                                    \end{feynman}
                                                \end{tikzpicture} \\ \hline
   $\gamma^*q\bar q$ & $ieQ_q\gamma^\mu\delta_{ij}$            & \begin{tikzpicture}[baseline=(current bounding box.center)]
                                                    \begin{feynman}
                                                        \vertex (a) at (0,0);
                                                        \vertex (b) at (0.8,0);
                                                        \vertex (c) at (1.5,0.5);
                                                        \vertex (d) at (1.5,-0.5);
                                                        \diagram* {
                                                        (a) -- [photon] (b),
                                                        (b) -- [fermion] (c),
                                                        (d) -- [fermion] (b)
                                                        };
                                                    \end{feynman}
                                                \end{tikzpicture} \\ \hline
  \end{tabular}
   
\label{tab:vertices}
\end{table}

\section{Passarino-Veltman Reduction}\label{subsec:PaVe}

At one-loop, amplitudes produce tensor integrals similar
to those in Eq.~\eqref{eq:PaVeSchematic} \cite{Passarino:1979pv}. With the massless\footnote{
Throughout this work the internal masses vanish and the physical external
partons satisfy $p_a^2=0$. The shifted momenta $r_i$, introduced in Eq.~\eqref{eq:paveDens},
entering propagators are combinations of external momenta and need not be lightlike.}
denominators defined as,
\begin{equation}\label{eq:paveDens}
    D_0 = \ell^2+i0,\quad D_i = (\ell + r_i)^2 + i0
\end{equation}
the tensor integral takes the form,
\begin{equation}
    I_N^{\mu_1...\mu_R} = \mu^{4-d}\int\frac{d^d\ell}{i\pi^{d/2}}\frac{\ell^{\mu_1}...\ell^{\mu_R}}{D_0D_1...D_{N-1}}
\end{equation}
where $R$ is the tensor rank and $r_i$ are the combinations of 
external momenta. 
The reduction is carried out in $d=4-2\epsilon$ dimensions, so that,
\begin{equation}
    g_{\mu\nu}g^{\mu\nu} = d.
\end{equation}

Passarino-Veltman reduction removes the Lorentz indices in the numerator 
$\ell^\mu\ell^\nu\ell^\rho...$ terms and expresses them 
in terms of external momenta $p_i^\mu$, the metric 
tensor $g_{\mu\nu}$ and scalar one-loop integrals, through 
Lorentz covariance \cite{Passarino:1979pv}.

For a rank-one $N$-point integral, the only available 
vectors are the independent external momenta, such that,
\begin{equation}\label{eq:paveRankOne}
    I_N^\mu = \sum_{i=1}^{N-1}p_i^\mu C_i
\end{equation}
where $C_i$ are scalar coefficient functions.
For rank-two, the relation becomes,
\begin{equation}
    I_N^{\mu\nu} = \sum_{i,j=1}^{N-1}p_i^\mu p_j^\nu C_{ij} + g^{\mu\nu}C_{00}
\end{equation}
and for rank-three,
\begin{equation}
    I_N^{\mu\nu\rho} = \sum_{i,j,k=1}^{N-1}p_i^\mu p_j^\nu p_k^\rho C_{ijk} + \sum_{i=1}^{N-1} \left(g^{\mu\nu}p_i^\rho + g^{\mu\rho}p_i^\nu + g^{\nu\rho}p_i^\mu\right)C_{00i}
\end{equation}
and so on for higher ranks.

The reduction mechanism works because scalar products 
involving the loop momentum can be rewritten using the
propagator denominators as,
\begin{equation}\label{eq:ellToDens}
    2\ell\cdot r_i = D_i - D_0 - r_i^2
\end{equation}
which lets us cancel denominator terms, through a method 
like the following example set at $i=1$,
\begin{equation}
    \frac{\ell\cdot r_1}{D_0D_1} = \frac12\left(\frac{1}{D_0}-\frac1{D_1}-\frac{r_1^2}{D_0D_1}\right)
\end{equation}
thereby rewriting loop-momentum scalar products in terms of denominators and 
scalar integrals.

While using this reduction scheme, we might also encounter Gram matrices, defined as,
\begin{equation}
    G_{ji} = 2p_j\cdot p_i.
\end{equation}
Contracting a general rank-one integral, as described in Eq.~\eqref{eq:paveRankOne},
with $p_j$ gives,
\begin{equation}
    p_j\cdot I_N = \sum_i(p_j\cdot p_i)C_i = \frac12\sum_iG_{ji}C_i.
\end{equation}
If we now define $R_j\equiv p_j\cdot I_N$, we are then able to compute the $C_i$
coefficients in the Lorentz decomposition by taking,
\begin{equation}\label{eq:gramInversion}
    G_{ji}C_i = 2R_j\Rightarrow C_i = 2\left(G^{-1}\right)_{ij}R_j.
\end{equation}

\subsection{Rank-One Bubble Example}

To illustrate the method, let us consider the rank-one bubble,
\begin{equation}
    B^\mu\!\left(p;0,0\right)=\mu^{4-d}\int\frac{d^d\ell}{i\pi^{d/2}}\frac{\ell^\mu}{D_0D_1}
\end{equation}
where the denominators follow Eq.~\eqref{eq:paveDens}. Lorentz covariance gives,
\begin{equation}
    B^\mu = p^\mu B_1
\end{equation}
Contracting with $p_\mu$ yields,
\begin{equation}
    p^2B_1 = p_\mu B^\mu = \mu^{4-d}\int\frac{d^d\ell}{i\pi^{d/2}}\frac{\ell\cdot p}{D_0D_1}.
\end{equation}
Using the relation in Eq.~\eqref{eq:ellToDens}, setting $r_1 = p$, we get,
\begin{equation}
    p^2B_1 = \frac12\left[A_0\!\left(0\right)-A_0\!\left(0\right)-p^2B_0\!\left(p^2;0,0\right)\right] = -\frac12p^2B_0\!\left(p^2;0,0\right)
\end{equation}
and therefore\footnote{The result shown in Eq.~\eqref{eq:bubbleResult} assumes generic 
kinematics with $p^2\neq 0$. Exceptional cases are obtained through an appropriate 
limit or alternative reduction.},
\begin{equation}\label{eq:bubbleResult}
    B^\mu = -\frac12p^\mu B_0\!\left(p^2;0,0\right).
\end{equation}
Here, $A_0$ and $B_0$ are the scalar one-loop integrals mentioned throughout. These 
two types can be defined as,
\begin{equation*}
    A_0\!\left(0\right) = \mu^{4-d}\int\frac{d^d\ell}{i\pi^{d/2}}\frac1{\ell^2+i0} = 0
\end{equation*}
\begin{equation}
    B_0\!\left(p^2;0,0\right) = \mu^{4-d}\int\frac{d^d\ell}{i\pi^{d/2}}\frac1{D_0D_1}.
\end{equation}
For massless propagators $A_0$ is scaleless and therefore vanishes in dimensional regularisation.

\subsection{Passarino-Veltman Scalar Functions}

As mentioned, only three of the MI considered are Passarino-Veltman scalar functions. Those are 
$B_0$, $C_0$ and $D_0$ and are the Passarino-Veltman two-, three- and four-point 
scalar integrals, respectively. The Passarino-Veltman method also features the $A_0$ (one-point)
scalar function, that does not feature in \textit{AntCalc}. It also accounts for 
scalar integrals with more than four external legs, as those can be further reduced 
in four dimensions to combinations of $D_0$ functions \cite{Melrose:1965kb,vanNeerven:1983vr}.

\subsubsection{The $B_0$ Function}\label{subsubsec:B0}

The $B_0$ function represents a scalar one-loop diagram with two external lines and 
two internal propagators. It is represented, in $d$ dimensions as,
\begin{equation}
    B_0\!\left(k^2;0,0\right) = \frac{(2\pi\mu)^{4-d}}{i\pi^2}\int d^dp\frac1{\left(p^2+i\varepsilon\right)\Big((p+k)^2+i\varepsilon\Big)}
\end{equation}
where $k$ is the external momentum and $\mu$ is the 't~Hooft normalisation scale.

It is also noteworthy that this integral, unlike higher-point ones, is UV divergent and 
generates the pole term $\Delta_B$, 
\begin{equation}
    \Delta_B = \frac2{4-d}-\gamma_E+\log(4\pi)
\end{equation}
where $\gamma_E$ is Euler's gamma constant.
This function may also be written as an integral over the Feynman parameter $x$. That 
yields the explicit analytical form \cite{Passarino:1979pv},
\begin{equation}
    \begin{split}
        B_0\!\left(k^2;0,0\right) &= \Delta_B - \int_0^1 dx\,\log\!\left[\frac{-x(1-x)k^2-i\varepsilon}{\mu^2}\right] + \mathcal O(d-4)\\
        &= \Delta_B + 2 - \log\!\left(\frac{-k^2-i\varepsilon}{\mu^2}\right) + \mathcal O(d-4).
    \end{split}
\end{equation}

\subsubsection{The $C_0$ Function}

The $C_0$ function represents a scalar loop with three external lines and three
internal propagators. It is a common feature in the vertex corrections of 
electroweak physics. It is written in $d$ dimensions as,
\begin{equation}
    \begin{split}
        C_0\!\left(k_1^2,k_2^2,(k_1+k_2)^2;0,0,0\right)=&\,\frac{(2\pi\mu)^{4-d}}{i\pi^2}\\ &\int d^dp\,\frac1{\Big(p^2+i\varepsilon\Big)\Big((p+k_1)^2+i\varepsilon\Big)\Big((p+k_1+k_2)^2+i\varepsilon\Big)}.
    \end{split}
\end{equation}

This integral is UV finite and may be rewritten as a double integral over the Feynman 
parameters $x$ and $y$ as,
\begin{equation}
    C_0\propto\int_0^1dx\int_0^{1-x}dy\,\left[ax^2+by^2+cxy+dx+ey+f\right]^{-1}
\end{equation}
where the coefficients $a, b, c, d, e$ and $f$ are linear combinations of the external 
momentum invariants \cite{tHooft:1978xw}.

\subsubsection{The $D_0$ Function}\label{subsubsec:D0}

The $D_0$ function describes a scalar loop with four external lines and four internal 
propagators. It is often called the scalar box integral. It is written as,
\begin{equation}
    D_0\!\left(k_1^2,k_2^2,k_3^2;k_4^2,(k_1+k_2)^2,(k_2+k_3)^2;0,0,0,0\right)=\frac{(2\pi\mu)^{4-d}}{i\pi^2}\int d^dp\,\frac1{D_0D_1D_2D_3}
\end{equation}
with denominator factors,
\begin{equation}
    \begin{split}
        &D_0 = p^2+i\varepsilon\\
        &D_i = \left(p+\sum_{a=1}^ik_a\right)^2+i\varepsilon,
    \end{split}
\end{equation}
with $k_4=-(k_1+k_2+k_3)$.

Evaluating $D_0$ is highly non-trivial as it 
features numerous physical cuts and branch points. The resulting analytical 
expression is expressed in terms of Spence functions (dilogarithms) \cite{tHooft:1978xw}.


\section{IBP Reduction Example}\label{sec:IBPRedEx}

The method behind IBP reduction, as discussed in Subsection~\ref{subsec:ibp} 
may be made easier to visualise with an example. Consider a one-loop 
bubble, akin to that studied via Passarino-Veltman reduction in Section~\ref{subsec:PaVe}. 
Take,
\begin{equation}
    \mathcal I(a_0,a_1) = \int d^d\ell\,\frac{1}{D_0^{a_0}D_1^{a_1}}
\end{equation}
where,
\begin{equation}
    D_0 = \ell^2-m_0^2,\quad D_1 = (\ell + p)^2 - m_1^2.
\end{equation}
We now choose $v^\mu = \ell^\mu$, and build and expand the relation from 
Eq.~\eqref{eq:derivativeToZero} as,
\begin{equation}
    \begin{split}
        0 &= \int d^d\ell\,\frac\partial{\partial\ell^\mu}\left(\frac{\ell^\mu}{D_0^{a_0}D_1^{a_1}}\right)\\
          &= d\mathcal I(a_0,a_1) - 2a_0\int d^d\ell\,\frac{\ell^2}{D_0^{a_0+1}D_1^{a_1}}-2a_1\int d^d\ell\,\frac{\ell\cdot (\ell+p)}{D_0^{a_0}D_1^{a_1+1}}.
    \end{split}
\end{equation}

We can now use,
\begin{gather}
    \ell^2 = D_0 + m_0^2\\
    2\ell\cdot (\ell + p) = D_0 + D_1 + m_0^2 + m_1^2 - p^2
\end{gather}
and finally get the explicit shifted-index identity,
\begin{equation}
    \begin{split}
        0 = \,&(d-2a_0-a_1)\mathcal I(a_0,a_1)\\
            &- 2a_0m_0^2\mathcal I(a_0+1,a_1)\\
            &- a_1\mathcal I(a_0-1,a_1+1)\\
            &-a_1\!\left(m_0^2+m_1^2-p^2\right)\mathcal I(a_0, a_1+1).
    \end{split}
\end{equation}

\section{Three-Particle Phase-Space to Cut Propagators}\label{sec:PhiToCuts}

Let us now consider the massless three-particle phase-space $\Phi_3$. 
It may be described as,
\begin{equation}
    d\Phi_3(q:p_1,p_2,p_3)\propto d^dp_1 d^dp_2 d^dp_3\,\delta^+\!\left(p_1^2\right)\,\delta^+\!\left(p_2^2\right)\,\delta^+\!\left(p_3^2\right)\,\delta^{(d)}(q-p_1-p_2-p_3).
\end{equation}
We start by using momentum conservation to eliminate $p_3$,
\begin{equation}
    p_3 = q - p_1 - p_2.
\end{equation}
This first step eliminates the rightmost $\delta$-function.
Then, using reverse unitarity, each of the remaining $\delta$-functions may be 
rewritten as cut propagators. Schematically, it gives,
\begin{equation}
    d\Phi_3\sim d^dp_1 d^dp_2 \frac1{[p_1^2]_c[p_2^2]_c[(q-p_1-p_2)^2]_c}
\end{equation}
with
\begin{equation}
    \left[\frac1{p^2}\right]_c\equiv \frac1{2\pi i}\left[\frac1{p^2-i0}-\frac1{p^2+i0}\right]_+.
\end{equation}

\section{Laporta Algorithm Implementation Example}\label{sec:laportaByHand}

Consider an arbitrary family defined as $\mathcal F[\{D_1, D_2, D_3\};d\Phi]$,
with the denominators chosen as $D_1 = s_{12}$, $D_2 = s_{13}$ and $D_3 = s_{23}$. 
We may take the integrals below, all members of this family, and organise them 
using their respective sector vectors as follows,
\begin{gather}
    I_1 = \int d\Phi\,\frac{s_{13}}{s_{12}}\Rightarrow I_1=I(1,-1,0),\\
    I_2 = \int d\Phi\,\frac2{s_{12}s_{13}s_{23}}\Rightarrow I_2=2I(1,1,1),\\
    I_3 = \int d\Phi\,\frac1{s_{12}^3}\Rightarrow I_3=I(3,0,0)
\end{gather}
where $d\Phi$ is the common phase-space measure.
Importantly, two integrals may share a sector if they contain the same denominators,
independently of those denominators' powers. Also, only positive indices denote 
denominator presence, meaning that integrals $I_1$ and $I_3$ in the example above 
are members of the same sector, $I_1, I_3\in (1,0,0)$. In our example, only two 
sectors appear: $(1,0,0)$ and $(1,1,1)$. These are our two target sectors. The 
three integrals are organised in terms of the $t$, $r$ and $s$ as follows,
\begin{equation}
    \begin{array}{c|ccc}
        & t & r & s\\
        \hline
        I(1,-1,0) & 1 & 0 & 1\\
        I(1,1,1)  & 3 & 0 & 0\\
        I(3,0,0)  & 1 & 2 & 0
    \end{array}
\end{equation}
then, by order of complexity, the algorithm orders the integrals as,
\begin{equation}
    I(1,1,1) \succ I(3,0,0) \succ I(1,-1,0).
\end{equation}

\section{The Massive Extension}\label{sec:massiveExtension}

The main body of this thesis focused primarily on the massless regime, since the
vast majority of the present work's chapters focus on that regime. This appendix
collects the supporting extension of the discussed methods into the massive regime.

The fundamental change is the massive on-shell condition; the remainder of this
section traces its consequences for the antenna construction and integration.
Since final-state particles now have mass, their squared momenta are not
set to zero, they are set to their squared mass.
\begin{equation}
    p_q^2 = p_{\bar q}^2 = 0\,\xrightarrow{\text{Massive Regime}}\, p_q^2=p_{\bar q}^2=m_q^2.
\end{equation}
where $p_i$ represents the momentum and $m_i$ represents the mass
of particle $i$.

The antenna subtraction formalism and the reduction architecture that was developed
and described throughout Chapter~\ref{ch:physicsbackground} are retained, however. What changes is the
intermediate steps and kinematics: the massive on-shell condition changes the
kinematics themselves and the unresolved limits, with soft radiation remaining
singular but heavy-quark collinear behaviour becoming quasi-collinear. Massive
integration is also different, as the characteristic phase-space scale in the massive
regime is not equal to its massless counterpart. A new scale of $m_q^2/q^2$ is also introduced and leads to
distinct cut-integral families and master integrals \cite{GehrmannDeRidder:2009fz}.

\subsection{Massive Kinematics and Infrared Structure}

In Subsection~\ref{subsec:infDivergences}, the Mandelstam invariants were introduced
as $s_{ij}=2p_i\cdot p_j$. This definition is retained in the massive regime, as in
Section~\ref{sec:massiveA30}, so that the squared invariant mass of a pair of partons
now also carries their masses,
\begin{equation}
    (p_i+p_j)^2 = s_{ij} + m_i^2 + m_j^2.
\end{equation}
Through momentum conservation, $q=\sum_i p_i$, the total invariant mass likewise becomes,
\begin{equation}
    q^2 = \sum_i m_i^2 + \sum_{i<j}s_{ij}.
\end{equation}

Soft radiation remains singular in the massive regime and therefore continues to
generate explicit infrared poles after integration. A non-zero mass regulates
collinear limits involving a massive parton, whereas collinear
limits between two massless partons remain singular.
If the two final-state collinear partons are massless, then the limit
follows the massless collinear limit from the massless regime. If the two partons
are massive, no strict collinear singularity exists. But if one of the collinear
partons is massive, consider that particle $i$ of mass $m_i$ and energy $E_i$, and the other
is massless, consider that particle $j$, such that $m_j=0$, then,
\begin{equation}
    p_i\cdot p_j = E_iE_j(1-\beta_i\cos\theta),\quad \beta_i=\sqrt{1-\frac{m_i^2}{E_i^2}}<1.
\end{equation}
As the angle between the partons $\theta$ tends towards zero, the mass prevents the
invariant from reaching the collinear singular surface,
\begin{equation}
    \lim_{\theta\to 0}p_i\cdot p_j = E_iE_j(1-\beta_i)\neq 0.
\end{equation}

This regulated collinear region still carries physically important information,
particularly in the joint high-energy and small-angle region,
$E_i>>m_i$ and $\theta<<1$,
\begin{equation}
    p_i\cdot p_j\sim\frac{E_iE_j}{2}\left(\theta^2+\frac{m_i^2}{E_i^2}\right)
\end{equation}
and consequently since,
\begin{equation}
    \int_0^{\theta_\text{max}^2} d\theta^2\,\frac{1}{p_i\cdot p_j}\sim\int_0^{\theta_\text{max}^2} d\theta^2\,\left(\theta^2+\frac{m_i^2}{E_i^2}\right)^{-1}\sim\ln\!\left(\frac{E_i^2}{m_i^2}\right)
\end{equation}
the massive-massless quasi-collinear region replaces the strict collinear pole by
a mass logarithm. The
quasi-collinear limit, then, is the limit in which both the mass $m_i$ and the opening angle
$\theta$ become small together. It is through this limit that the massive antennae
recover the massless splitting behaviour, as $m_i\to 0$ \cite{GehrmannDeRidder:2009fz}.

Mass does not regulate UV divergences, as $\ell^2>>m_i^2,p_i^2$. It is noteworthy,
however, that in a theory with massive quarks, the coupling renormalisation of
Eq.~\eqref{eq:alphaRescaling} must be supplemented by the renormalisation of the quark mass
and quark field,
\begin{equation}
    m_{q,0} = Z_m m_q,\quad \psi_{q,0} = \sqrt{Z_2^q}\psi_q.
\end{equation}

\subsection{Massive Phase-Space Integration and Reduction Methods}

The $n$-particle phase-space measure of Eq.~\eqref{eq:nParticlePS} already carries the
final-state masses through its on-shell conditions $\delta^+(p_i^2-m_i^2)$. The massless
regime of the main text sets $m_i=0$, whereas in the massive regime these conditions
impose $p_i^2=m_i^2$, with $p_i^0>0$, for the heavy quarks.

The phase-space factorisations of Eqs.~\eqref{eq:nloPhaseSpaceFactor} and
\eqref{eq:nnloPhaseSpaceFactor} also still hold, with the antenna phase space evaluated
at $q=\tilde p_I+\tilde p_K$ (or $q=\tilde p_I+\tilde p_L$), which expresses momentum
conservation across the mapping. The only change is that the momentum mapping must now
also preserve the mass-shell conditions of the hard radiators, $\tilde p_I^2 = m_i^2$ and
$\tilde p_K^2 = m_k^2$.

Antenna integration is computed in the same manner as in the massless regime,
\begin{equation}
    \mathcal X_n^l\sim\int d\Phi_X\,\left(\prod_{r=1}^l d^d\ell_r\right)X_n^l(\{p\},\{\ell\})\sim\left(q^2\right)^{\lambda+\kappa\epsilon}f\!\left(\left\{\frac{m^2}{q^2}\right\},\epsilon\right).
\end{equation}

Passarino-Veltman and IBP reduction remain applicable in the massive regime, following the same
principles introduced for the massless case.
For IBP reduction, however, the massive case is significantly harder analytically:
the phase-space is non-trivial and has mass-dependent thresholds; the master
integrals are generally multi-scale, and can involve
logarithms, polylogarithms and more complicated transcendental functions;
and lastly, boundary conditions are less direct, although common choices are
the massless limit at $m\to0$ and the threshold limit at $\beta\to0$ \cite{GehrmannDeRidder:2009fz}.

\section{Master Integral Substitution Arrays}\label{sec:MISubsArrays}

As discussed in Section~\ref{sec:laportaByHand}, \textit{AntCalc}'s IBP reduction,
handled internally by \textit{LiteRed2}, organises each antenna family's master
integrals as members of that family's own basis, addressed through \textit{LiteRed2}'s
$j[\text{basis},...]$ notation rather than the named masters used throughout
Chapter~\ref{ch:workedexample}. At the master combination stage for each antenna, the
raw $j[...]$-form result is substituted through a rule list specific to that family's
basis, translating it into the named masters. The rule lists used are collected below,
one per antenna family.

\subsection{The $A_2^2$ Set}

\begin{mdframed}[backgroundcolor=gray!10, linewidth=0pt]
\texttt{In[1]:= A22MISubs = \{\\
   j[A22TwoLoopTreeBasis1, 0, 0, 0, 1, 1, 0, 1] -> A3,\\
   j[A22TwoLoopTreeBasis1, 1, 0, 1, 1, 1, 0, 0] -> A4,\\
   j[A22TwoLoopTreeBasis4, 1, 0, 1, 0, 0, 1, 0] -> A3,\\
   j[A22TwoLoopTreeBasis4, 1, 0, 0, 1, 0, 1, 1] -> A4,\\
   j[A22TwoLoopTreeBasis5, 0, 0, 1, 0, 0, 1, 1] -> A3,\\
   j[A22TwoLoopTreeBasis6, 0, 0, 1, 1, 0, 0, 1] -> A3,\\
   j[A22TwoLoopTreeBasis6, 1, 0, 1, 1, 0, 1, 0] -> A4,\\
   j[A22TwoLoopTreeBasis7, 0, 0, 1, 0, 1, 0, 1] -> A3,\\
   j[A22TwoLoopTreeBasis7, 1, 0, 1, 0, 1, 1, 0] -> A4,\\
   j[A22TwoLoopTreeBasis7, 1, 1, 0, 1, 0, 1, 0] -> A22LO,\\
   j[A22TwoLoopTreeBasis8, 0, 1, 0, 0, 0, 1, 1] -> A3,\\
   j[A22TwoLoopTreeBasis8, 1, 0, 1, 1, 1, 0, 0] -> A4,\\
   j[A22TwoLoopTreeBasis8, 1, 1, 1, 1, 1, 1, 0] -> A6,\\
   j[A22OneLoopSelfBasis, 1, 0, 1, 1, 0, 1, 0] -> A22LO\\
   \};}
\end{mdframed}

\subsection{The $A_3^1$ Set}

\begin{mdframed}[backgroundcolor=gray!10, linewidth=0pt]
\texttt{In[1]:= A31MISubs = \{\\
   j[A31Basis10, 1, 1, 1, 0, 0, 1, 0, 1, 0] -> V5b,\\
   j[A31Basis12, 1, 1, 1, 0, 0, 1, 0, 1, 0] -> V5b,\\
   j[A31Basis2, 1, 1, 1, 0, 0, 1, 0, 1, 0] -> V5b,\\
   j[A31Basis6, 1, 1, 1, 0, 1, 0, 0, 1, 0] -> V5a,\\
   j[A31Basis6, 1, 1, 1, 0, 0, 1, 0, 1, 0] -> V5b,\\
   j[A31Basis9, 1, 1, 1, 1, 0, 0, 0, 1, 0] -> V5a,\\
   j[A31Basis9, 1, 1, 1, 0, 1, 0, 0, 1, 0] -> V5b,\\
   j[A31Basis3, 1, 1, 1, 0, 1, 0, 0, 1, 0] -> V5a,\\
   j[A31Basis3, 1, 1, 1, 0, 0, 0, 1, 1, 0] -> V5b,\\
   j[A31Basis4, 1, 1, 1, 0, 1, 0, 0, 1, 0] -> V5a,\\
   j[A31Basis4, 1, 1, 1, 0, 0, 1, 0, 1, 0] -> V5b,\\
   j[A31Basis4, 1, 1, 1, 0, 1, 1, 1, 1, 1] -> 2 V8,\\
   j[A31Basis7, 1, 1, 1, 1, 0, 0, 0, 1, 0] -> V5a,\\
   j[A31Basis7, 1, 1, 1, 0, 1, 0, 0, 1, 0] -> V5b,\\
   j[A31Basis7, 1, 1, 1, 1, 1, 1, 1, 1, 0] -> V8,\\
   j[A31Basis8, 1, 1, 1, 1, 0, 0, 0, 1, 0] -> V5a,\\
   j[A31Basis8, 1, 1, 1, 0, 1, 0, 0, 1, 0] -> V5b\\
   \};}
\end{mdframed}

\subsection{The Four-Parton (X40) Set}

\begin{mdframed}[backgroundcolor=gray!10, linewidth=0pt]
\texttt{In[1]:= X40MISubs = \{\\
   j[chainBasis1342, 1, 1, 1, 1, 0, 0, 0, 0, 0] -> R4,\\
   j[chainBasis1342, 1, 1, 1, 1, 0, 0, 0, 1, 1] -> R6,\\
   j[chainBasis1432, 1, 1, 1, 1, 0, 0, 0, 0, 0] -> R4,\\
   j[chainBasis1432, 1, 1, 1, 1, 0, 0, 0, 1, 1] -> R6,\\
   j[boxBasis1324, 1, 1, 1, 1, 0, 0, 0, 0, 0] -> R4,\\
   j[boxBasis1324, 1, 1, 1, 1, 1, 1, 1, 1, 0] -> R8a,\\
   j[chainBasis1324, 1, 1, 1, 1, 0, 0, 0, 0, 0] -> R4,\\
   j[chainBasis1423, 1, 1, 1, 1, 0, 0, 0, 0, 0] -> R4,\\
   j[hybridBasis1234, 1, 1, 1, 1, 0, 0, 0, 0, 0] -> R4,\\
   j[hybridBasis1234, 1, 1, 1, 1, 0, 0, 0, 1, 1] -> R6,\\
   j[hybridBasis1234, 1, 1, 1, 1, 1, 0, 1, 1, 1] -> R8b,\\
   j[hybridBasis1243, 1, 1, 1, 1, 0, 0, 0, 0, 0] -> R4,\\
   j[hybridBasis1243, 1, 1, 1, 1, 0, 0, 0, 1, 1] -> R6,\\
   j[hybridBasis1243, 1, 1, 1, 1, 1, 0, 1, 1, 1] -> R8b,\\
   j[chainBasis1234, 1, 1, 1, 1, 0, 0, 0, 0, 0] -> R4,\\
   j[chainBasis1234, 1, 1, 1, 1, 0, 0, 0, 1, 1] -> R6,\\
   j[chainBasis2134, 1, 1, 1, 1, 0, 0, 0, 0, 0] -> R4,\\
   j[chainBasis2134, 1, 1, 1, 1, 0, 0, 0, 1, 1] -> R6,\\
   j[chainBasis2314, 1, 1, 1, 1, 0, 0, 0, 0, 0] -> R4,\\
   j[chainBasis2314, 1, 1, 1, 1, 0, 0, 0, 1, 1] -> R6,\\
   j[hybridBasis2431, 1, 1, 1, 1, 0, 0, 0, 0, 0] -> R4,\\
   j[hybridBasis2431, 1, 1, 1, 1, 0, 0, 0, 1, 1] -> R6,\\
   j[hybridBasis2431, 1, 1, 1, 1, 1, 0, 1, 1, 1] -> R8b\\
   \};}
\end{mdframed}

\section{Local Validation of Additional Routes}\label{sec:localValidationExtra}

Much like the $A_3^0$ massless unintegrated antenna was locally validated by computing its local soft and 
(single) collinear limits, so have other antenna functions been validated in that way. In this section 
we present those validation methods and results.

\subsection{$A_4^0$}\label{subsec:localValidationA40}

The leading $A_4^0(1_q,3_g,4_g,2_{\bar q})$ massless antenna component has five relevant limits: two soft limits,
and three collinear limits, for $1_q||3_g$, $3_g||4_g$, and $2_{\bar q}||4_g$. The two $qg$ collinear limits, $1_q||3_g$ and $2_{\bar q}||4_g$,
and the two soft-$g$ limits are symmetric under the reversal symmetry of the colour 
ordering, as $1\leftrightarrow 2$ and $3\leftrightarrow4$.


\subsubsection{Soft Limit and the Eikonal Factor}

In this limit we set $p_g\to 0$. As mentioned, the limit is symmetric for both gluons. We will be focusing on the 
soft-$3_g$ here.

The limit is thus computed as,
\begin{equation}
    \lim_{\lambda\to 0}\lambda^2 A_4^0\Big|_{s_{13},s_{23},s_{34}\to \lambda s_{13},\lambda s_{23},\lambda s_{34}} = \mathcal S_{14}^{(3)}A_3^0(1_q,4_g,2_{\bar q}),\quad \mathcal S_{14}^{(3)}=\frac{2s_{14}}{s_{13}s_{34}}.
\end{equation}

This limit is computed using \textit{AntCalc} as follows \cite{Gehrmann-DeRidder:2005btv}.

\begin{mdframed}[backgroundcolor=gray!10, linewidth=0pt]
\begin{lstlisting}
In[1]:= A40 = BuildAntenna[A,4,0,Component->Leading];

        soft3Rules = {s13 -> lambda s13, s23 -> lambda s23,
            s34 -> lambda s34, s134 -> s14, s234 -> s24, q2 -> s124};
        SeriesCoefficient[A40 /. soft3Rules, {lambda, 0, -2}] // Simplify
\end{lstlisting}
\texttt{Out[1]:=}
    \begin{equation*}
        \frac{4s_{12}^2+4s_{12}(s_{14}+s_{24})-2(\epsilon-1)s_{14}^2-4\epsilon s_{14}s_{24}-2(\epsilon-1)s_{24}^2}{s_{124}s_{13}s_{24}s_{34}}
    \end{equation*}
\end{mdframed}
Since the result is larger than a simple factor, unlike for $A_3^0$, it is compared with the right-hand side of the
equation above after the calculation. Using Eq.~\eqref{eq:A30} for $A_3^0(1_q,4_g,2_{\bar q})$, with $s_{13}\to s_{14}$,
$s_{23}\to s_{24}$ and $q^2\to s_{124}$, that right-hand side is
\begin{equation}
    \mathcal S_{14}^{(3)}A_3^0(1_q,4_g,2_{\bar q})=\frac{2s_{14}}{s_{13}s_{34}}\,
    \frac{2s_{12}s_{124}+s_{14}^2+s_{24}^2-\epsilon(s_{14}+s_{24})^2}{s_{124}s_{14}s_{24}},
\end{equation}
which coincides with the result above once $s_{124}=s_{12}+s_{14}+s_{24}$ is used.


\subsubsection{Collinear $1_q||3_g$ Limit}

For the $1_q||3_g$ limit, the quark and gluon pair recombine into a parent quark as,
\begin{equation}
    p_1=zp_{13},\quad p_3=(1-z)p_{13}
\end{equation}
where $z$ is the momentum fraction carried by the original quark. This step makes $s_{13}$ vanish as, $s_{13}\to 0$.

We now introduce the Mandelstam notation $s_{(13)j}$ to represent the invariant computed from the momenta of the 
parent quark and some other $j$ particle. The remaining invariants become,
\begin{equation}
    \begin{gathered}
        s_{12}\to zs_{(13)2},\quad s_{23}\to(1-z)s_{(13)2},\\
        s_{14}\to zs_{(13)4},\quad s_{34}\to(1-z)s_{(13)4},\\
        s_{24}\to s_{24}.
    \end{gathered}
\end{equation} 
We also define the reduced configuration of the three-parton tree-level antenna as,
\begin{equation}
    A_3^0\Big((13)_q,4_g,2_{\bar q}\Big).
\end{equation}

At leading power, the four-parton antenna factorises as,
\begin{equation}
    \lim_{s_{13}\to 0} s_{13}A_4^0(1_q,3_g,4_g,2_{\bar q})=\frac{1}{s_{13}}P_{qg\to q}(z,\epsilon)A_3^0\Big((13)_q,4_g,2_{\bar q}\Big),\quad P_{qg\to q}(z,\epsilon)=\frac{1+z^2}{1-z}-\epsilon(1-z).
\end{equation}
Here $P_{qg\to q}(z,\epsilon)$ is the $\mathcal O(\epsilon)$-extended version of the splitting
function used for $A_3^0$'s collinear limit in Subsection~\ref{subsec:a30LocalValidation}; the
$-\epsilon(1-z)$ term is the known $d$-dimensional extension from the literature, quoted here
since $A_4^0$ is an NNLO antenna where it no longer drops out \cite{Gehrmann-DeRidder:2005btv}.
The check below is performed at $\epsilon=0$, and therefore verifies the $(1+z^2)/(1-z)$ piece only.

\begin{mdframed}[backgroundcolor=gray!10, linewidth=0pt]
\begin{lstlisting}
In[1]:= A40 = BuildAntenna[A,4,0,Component->Leading];

        col13Rules = {s13 -> lambda, s12 -> z sP2, s23 -> (1-z) sP2,
            s14 -> z sP4, s34 -> (1-z) sP4, s134 -> sP4,
            s234 -> (1-z) sP2 + s24 + (1-z) sP4, q2 -> sP2 + sP4 + s24,
            Epsilon -> 0};
        SeriesCoefficient[Together[A40 /. col13Rules],
            {lambda, 0, -1}] // FullSimplify
\end{lstlisting}
\texttt{Out[1]:=}
    \begin{equation*}
        -\frac{(1+z^2)\left(s_{24}^2+2s_{24}s_{P2}+2s_{P2}^2+2s_{P2}s_{P4}+s_{P4}^2\right)}{s_{24}s_{P4}(z-1)(s_{24}+s_{P2}+s_{P4})}
    \end{equation*}
\end{mdframed}
This is the $\epsilon=0$ kernel $(1+z^2)/(1-z)$ times $A_3^0\big((13)_q,4_g,2_{\bar q}\big)$: with $s_{12}\to s_{P2}$,
$s_{13}\to s_{P4}$, $s_{23}\to s_{24}$ and $q^2\to s_{P2}+s_{P4}+s_{24}$, Eq.~\eqref{eq:A30} gives
\begin{equation*}
    A_3^0=\frac{s_{24}^2+2s_{24}s_{P2}+2s_{P2}^2+2s_{P2}s_{P4}+s_{P4}^2}{s_{24}s_{P4}(s_{24}+s_{P2}+s_{P4})}.
\end{equation*}


\subsubsection{Collinear $3_g||4_g$ Limit}

For the $3_g||4_g$ limit, the gluon-gluon pair recombine into a parent gluon as,
\begin{equation}
    p_3=zp_{34},\quad p_4=(1-z)p_{34}.
\end{equation}
This step makes $s_{34}$ vanish.

The remaining invariants become,
\begin{equation}
    \begin{gathered}
        s_{13}\to zs_{1(34)},\quad s_{14}\to(1-z)s_{1(34)},\\
        s_{23}\to zs_{(34)2},\quad s_{24}\to(1-z)s_{(34)2},\\
        s_{12}\to s_{12}.
    \end{gathered}
\end{equation} 
The reduced configuration of the three-parton tree-level antenna also becomes,
\begin{equation}
    A_3^0\Big(1_q,(34)_g,2_{\bar q}\Big).
\end{equation}

The parent parton here is now a gluon, not a quark. Hence, the local factorisation
formula involves the spin-correlated $g\to gg$ splitting tensor:
\begin{equation}
    A_4^0(1_q,3_g,4_g,2_{\bar q})\xrightarrow{3_g||4_g}\frac1{s_{34}}\hat P_{gg\to G}^{\mu\nu}(z,k_\perp)\big(A_3^0\big)_{\mu\nu}\Big(1_q,(34)_g,2_{\bar q}\Big).
\end{equation}
Here, $\big(A_3^0\big)_{\mu\nu}$ is the reduced three-parton antenna with the parent-gluon polarisation
indices left uncontracted. Equivalently, in terms of scalar antenna functions, the
unaveraged collinear limit contains an angular contribution,
\begin{equation}
    A_4^0(1_q,3_g,4_g,2_{\bar q})\xrightarrow{3_g||4_g}\frac1{s_{34}}P_{gg\to G}(z)A_3^0\Big(1_q,(34)_g,2_{\bar q}\Big)+\mathrm{ang.},
\end{equation}
where $\mathrm{ang.}$ denotes a spin-correlation contribution. This term
vanishes after integration over the unresolved azimuthal angle. Consequently,
\begin{equation}
    \lim_{s_{34}\to0}s_{34}\left\langle A_4^0(1_q,3_g,4_g,2_{\bar q})\right\rangle_\phi = P_{gg\to G}(z)A_3^0\Big(1_q,(34)_g,2_{\bar q}\Big),
\end{equation}
where
\begin{equation}
    P_{gg\to G}(z)=2\left[\frac{z}{1-z}+\frac{1-z}{z}+z(1-z)\right].
\end{equation}
The direct invariant-space limit produces a non-vanishing residual after
subtracting the spin-averaged contribution, consistently with the angular term
in the local gluon--gluon collinear factorisation \cite{Gehrmann-DeRidder:2005btv}.

\begin{mdframed}[backgroundcolor=gray!10, linewidth=0pt]
\begin{lstlisting}
In[1]:= A40 = BuildAntenna[A,4,0,Component->Leading];

        col34Rules = {s34 -> lambda, s13 -> z s1P, s14 -> (1-z) s1P,
            s23 -> z sP2, s24 -> (1-z) sP2, s134 -> s1P + lambda,
            s234 -> sP2 + lambda, q2 -> s12 + s1P + sP2 + lambda,
            Epsilon -> 0};
        SeriesCoefficient[Together[A40 /. col34Rules],
            {lambda, 0, -1}] // FullSimplify
\end{lstlisting}
\texttt{Out[1]:=}
    \begin{equation*}
        \frac{-4\left(1+2z(z-1)\right)s_{12}\left(s_{12}+s_{1P}+s_{P2}\right)-2\left(s_{1P}^2+s_{P2}^2\right)\left(1+z(z-1)\right)^2}{s_{1P}s_{P2}\,z(z-1)\left(s_{12}+s_{1P}+s_{P2}\right)}
    \end{equation*}
\end{mdframed}
Subtracting $P_{gg\to G}(z)A_3^0\big(1_q,(34)_g,2_{\bar q}\big)$, the spin-averaged term of the equation above, from this
result leaves exactly the angular term.

Note that the $\mathrm{ang.}$ quantity was found to be,
\begin{equation}
    \mathrm{ang.} = \frac{4s_{12}z(z-1)}{s_{1(34)}s_{(34)2}}
\end{equation}
but this value is not to be understood as a general result.

\subsection{$\tilde A_4^0$}\label{subsec:localValidationA40Tilde}

The subleading $\tilde A_4^0(1_q,3_g,4_g,2_{\bar q})$ massless antenna component has the same five
kinematic regions as the leading term, but their singular structure differs. As seen in
Subsection~\ref{subsec:antFamilies}, $\tilde A_4^0$ is obtained from the squared sum of both colour
orderings, $\left|\mathcal M_4^0(1_q,3_g,4_g,2_{\bar q})+\mathcal M_4^0(1_q,4_g,3_g,2_{\bar q})\right|^2$, so its
gluons are not colour ordered and it must be symmetric under the exchange of the two gluon labels,
$3\leftrightarrow4$. The check below confirms this: once the physical relations for $s_{134}$ and
$s_{234}$ are imposed, the difference between $\tilde A_4^0$ and its gluon-exchanged copy vanishes, so
the antenna is symmetric in the gluon indices.

\begin{mdframed}[backgroundcolor=gray!10, linewidth=0pt]
\begin{lstlisting}
In[1]:= tA40 = BuildAntenna[A,4,0,Component->Subleading];

        physicalRelations = {s134 -> s13 + s14 + s34,
            s234 -> s23 + s24 + s34};

        tA40Physical = Together[tA40 /. physicalRelations];

        gluonExchangeRules = {s13 -> s14, s14 -> s13, s23 -> s24,
            s24 -> s23};

        FullSimplify[tA40Physical - (tA40Physical /. gluonExchangeRules)]
\end{lstlisting}
\texttt{Out[1]:=}
    \begin{equation*}
        0
    \end{equation*}
\end{mdframed}

All other limits carry the same method as the ones for $A_4^0$.

\subsubsection{Soft Limit and the Eikonal Factor}

This limit is also symmetric for the $4_g$ gluon. It also generates the eikonal factor \cite{Gehrmann-DeRidder:2005btv}.

\begin{mdframed}[backgroundcolor=gray!10, linewidth=0pt]
\begin{lstlisting}
In[1]:= tA40 = BuildAntenna[A,4,0,Component->Subleading];

        soft3Rules = {s13 -> lambda s13, s23 -> lambda s23,
            s34 -> lambda s34, s134 -> s14, s234 -> s24, q2 -> s124};
        SeriesCoefficient[tA40 /. soft3Rules, {lambda, 0, -2}] // Simplify
\end{lstlisting}
\texttt{Out[1]:=}
    \begin{equation*}
        \frac{2s_{12}\left[2s_{12}(s_{12}+s_{14}+s_{24})+s_{14}^2+s_{24}^2-\epsilon(s_{14}+s_{24})^2\right]}{s_{124}s_{13}s_{14}s_{23}s_{24}}
    \end{equation*}
\end{mdframed}
The reduced antenna keeps its $\epsilon$ dependence. With Eq.~\eqref{eq:A30} for $A_3^0(1_q,4_g,2_{\bar q})$ as above, the
result equals, for general $\epsilon$,
\begin{equation}
    \frac{2s_{12}}{s_{13}s_{23}}A_3^0(1_q,4_g,2_{\bar q})=\frac{2s_{12}}{s_{13}s_{23}}\,
    \frac{2s_{12}s_{124}+s_{14}^2+s_{24}^2-\epsilon(s_{14}+s_{24})^2}{s_{124}s_{14}s_{24}},
\end{equation}
that is, the eikonal factor $\mathcal S_{12}^{(3)}$ between the quark and the antiquark times the reduced antenna.


\subsubsection{Collinear $1_q||3_g$ Limit}

The $1_q||3_g$ limit also behaves in the same way as for the leading term, again checked here at
$\epsilon=0$ against the $(1+z^2)/(1-z)$ kernel \cite{Gehrmann-DeRidder:2005btv}.

\begin{mdframed}[backgroundcolor=gray!10, linewidth=0pt]
\begin{lstlisting}
In[1]:= tA40 = BuildAntenna[A,4,0,Component->Subleading];

        col13Rules = {s13 -> lambda, s12 -> z sP2, s23 -> (1-z) sP2,
            s14 -> z sP4, s34 -> (1-z) sP4, s134 -> sP4,
            s234 -> (1-z) sP2 + s24 + (1-z) sP4, q2 -> sP2 + sP4 + s24,
            Epsilon -> 0};
        SeriesCoefficient[Together[tA40 /. col13Rules],
            {lambda, 0, -1}] // FullSimplify
\end{lstlisting}
\texttt{Out[1]:=}
    \begin{equation*}
        -\frac{(1+z^2)\left(s_{24}^2+2s_{24}s_{P2}+2s_{P2}^2+2s_{P2}s_{P4}+s_{P4}^2\right)}{s_{24}s_{P4}(z-1)(s_{24}+s_{P2}+s_{P4})}
    \end{equation*}
\end{mdframed}
The result is identical to that of the leading-colour component above.


\subsubsection{Collinear $3_g||4_g$ Limit}

The $3_g||4_g$ limit is the one that differs significantly from its analogue for $A_4^0$. Using the
same parametrisation as for $A_4^0$, the expansion of $\tilde A_4^0$ in $s_{34}$ starts at order
$s_{34}^0$: the coefficient of $1/s_{34}$, computed below for general $\epsilon$, vanishes identically,
so $\tilde A_4^0$ has no $3_g||4_g$ collinear singularity. The gluon-exchange symmetry shown above is
not, on its own, the reason for this, since the sum of the two colour-ordered $A_4^0$ terms is equally
symmetric and is singular. The reason is that the only diagram carrying the $1/s_{34}$ propagator is
the one with the three-gluon vertex, and, as the colour-ordered three-gluon vertex is antisymmetric
in its gluon legs, this diagram enters $\mathcal M_4^0(1_q,3_g,4_g,2_{\bar q})$ and
$\mathcal M_4^0(1_q,4_g,3_g,2_{\bar q})$ with opposite signs. It therefore cancels in their sum, in
which the two gluons couple to the quark line like photons \cite{Gehrmann-DeRidder:2005btv}.

\begin{mdframed}[backgroundcolor=gray!10, linewidth=0pt]
\begin{lstlisting}
In[1]:= tA40 = BuildAntenna[A,4,0,Component->Subleading];

        col34Rules = {s34 -> lambda, s13 -> z s1P, s14 -> (1-z) s1P,
            s23 -> z sP2, s24 -> (1-z) sP2, s134 -> s1P + lambda,
            s234 -> sP2 + lambda, q2 -> s12 + s1P + sP2 + lambda};

        FullSimplify[SeriesCoefficient[Together[tA40 /. col34Rules],
            {lambda, 0, -1}]]
\end{lstlisting}
\texttt{Out[1]:=}
    \begin{equation*}
        0
    \end{equation*}
\end{mdframed}

\subsection{$B_4^0$}\label{subsec:localValidationB40}

The $B_4^0$ antenna is the ordinary, non-identical-flavour four-quark antenna of the process $\gamma^*\to qq^\prime\bar q^\prime\bar q$.
This antenna has three non-vanishing limits, the single collinear limit $3_{q^\prime}||4_{\bar q^\prime}$, the double soft limit
$3_{q^\prime},4_{\bar q^\prime}\to 0$, and the triple-collinear limit $1_q||3_{q^\prime}||4_{\bar q^\prime}$.

\subsubsection{Single Collinear Limit}

This limit is computed similarly to $A_4^0$'s $3_g||4_g$ collinear limit. We find that,
\begin{equation}
    B_4^0(1_q,3_{q^\prime},4_{\bar q^\prime},2_{\bar q})\xrightarrow{3_{q^\prime}||4_{\bar q^\prime}}\frac1{s_{34}}P_{q\bar q\to g}^{\mu\nu}(z,k_\perp)\left(A_3^0\right)_{\mu\nu}\Big(1,(34),2\Big)
\end{equation}
which after azimuthal averaging becomes,
\begin{equation}
    \left\langle B_4^0\right\rangle_\phi \xrightarrow{3_{q^\prime}||4_{\bar q^\prime}}\frac1{s_{34}} P_{q\bar q\to g}(z)A_3^0\Big(1,(34),2\Big),\quad P_{q\bar q\to g}(z) = \frac{z^2+(1-z)^2-\epsilon}{1-\epsilon}.
\end{equation}
The check below keeps $\epsilon$ symbolic throughout, directly validating the full $\epsilon$-dependent
kernel $P_{q\bar q\to g}(z)$ above, rather than only its $\epsilon=0$ piece \cite{Gehrmann-DeRidder:2005btv}.

\begin{mdframed}[backgroundcolor=gray!10, linewidth=0pt]
\begin{lstlisting}
In[1]:= B40 = BuildAntenna[B,4,0];

        col34Rules = {s34 -> lambda, s13 -> z s1P, s14 -> (1-z) s1P,
            s23 -> z sP2, s24 -> (1-z) sP2, s134 -> s1P, s234 -> sP2,
            q2 -> s12 + s1P + sP2};
        SeriesCoefficient[Together[B40 /. col34Rules],
            {lambda, 0, -1}] // FullSimplify
\end{lstlisting}
\texttt{Out[1]:=}
    \begin{equation*}
        \frac{2(\epsilon-1)s_{12}\left(s_{12}+s_{1P}+s_{P2}\right)-\left(\epsilon-1-2z(z-1)\right)\left[(\epsilon-1)\left(s_{1P}^2+s_{P2}^2\right)+2\epsilon s_{1P}s_{P2}\right]}{(\epsilon-1)s_{1P}s_{P2}\left(s_{12}+s_{1P}+s_{P2}\right)}
    \end{equation*}
\end{mdframed}
Subtracting $P_{q\bar q\to g}(z)A_3^0\big(1_q,(34)_g,2_{\bar q}\big)$ from this result leaves exactly the
spin-correlation term $4s_{12}z(z-1)/\big[(\epsilon-1)s_{1P}s_{P2}\big]$, which vanishes after azimuthal averaging.


\subsubsection{Double-Soft Limit}

A soft quark-antiquark pair is singular when both particles become soft simultaneously. To 
compute this limit we start scalling,
\begin{equation}
    \begin{gathered}
        p_3\to \lambda p_3,\quad p_4\to \lambda p_4,\\
        s_{i3},s_{i4}\sim\lambda,\quad s_{34}\sim\lambda^2.
    \end{gathered}
\end{equation}

The antenna has a leading $\lambda^{-4}$ behaviour such that \cite{Gehrmann-DeRidder:2005btv},
\begin{equation}
    \begin{split}
    \lim_{\lambda\to0}\lambda^4 B_4^0(1,\lambda3,\lambda4,2)&=S_{12}(3,4),\\
    S_{12}(3,4) &= \frac{2\left(s_{12}s_{34}-s_{13}s_{24}-s_{14}s_{23}\right)}{s_{34}^2(s_{13}+s_{14})(s_{23}+s_{24})}+\frac{2}{s_{34}^2}\left[\frac{s_{13}s_{14}}{(s_{13}+s_{14})^2}+\frac{s_{23}s_{24}}{(s_{23}+s_{24})^2}\right]
    \end{split}
\end{equation}

\begin{mdframed}[backgroundcolor=gray!10, linewidth=0pt]
\begin{lstlisting}
In[1]:= B40 = BuildAntenna[B,4,0];

        doubleSoftRules = {s13 -> lambda s13, s14 -> lambda s14,
            s23 -> lambda s23, s24 -> lambda s24, s34 -> lambda^2 s34,
            s134 -> lambda (s13+s14) + lambda^2 s34,
            s234 -> lambda (s23+s24) + lambda^2 s34,
            q2 -> s12 + lambda (s13+s14+s23+s24) + lambda^2 s34,
            Epsilon -> 0};
        SeriesCoefficient[Together[B40 /. doubleSoftRules],
            {lambda, 0, -4}] // FullSimplify
\end{lstlisting}
\texttt{Out[1]:=}
    \begin{equation*}
        \frac{2s_{12}s_{34}(s_{13}+s_{14})(s_{23}+s_{24})-2\left(s_{14}s_{23}-s_{13}s_{24}\right)^2}{s_{34}^2(s_{13}+s_{14})^2(s_{23}+s_{24})^2}
    \end{equation*}
\end{mdframed}
This is the double-soft factor $\mathcal S_{12}(3,4)$ of the equation above, written over a common denominator.


\subsubsection{Triple-Collinear Limit}

For the $1_q||3_{q^\prime}||4_{\bar q^\prime}$ limit, all three partons become collinear together. We
introduce the rescaled invariants $r_{13}$, $r_{14}$, $r_{34}$, of order unity as $\lambda\to0$, via
\begin{equation}
    s_{13}\to\lambda r_{13},\quad s_{14}\to\lambda r_{14},\quad s_{34}\to\lambda r_{34},
\end{equation}
and the momentum fractions $w$, $x$, and $1-w-x$ carried by partons $1$, $3$, and $4$ respectively
within the recombined cluster, via
\begin{equation}
    s_{12}\to w\,s_{P2},\quad s_{23}\to x\,s_{P2},\quad s_{24}\to(1-w-x)\,s_{P2}.
\end{equation}

At leading power, the four-parton antenna factorises onto the trivial LO antenna, $A_2^0\equiv1$, as,
\begin{equation}
    \lim_{\lambda\to0}\lambda^2 B_4^0(1_q,3_{q^\prime},4_{\bar q^\prime},2_{\bar q}) = P_\mathrm{NonIdent}^{(134)}(w,x;r_{13},r_{14},r_{34})\,A_2^0(1,2),
\end{equation}
where the non-identical-flavour triple-collinear splitting function is \cite{Gehrmann-DeRidder:2005btv},
\begin{equation}
    \begin{split}
        P_\mathrm{NonIdent}^{(134)}(w,x;r_{13},r_{14},r_{34}) =\;&
            -\frac{(1-\epsilon)+\dfrac{2r_{13}}{r_{34}}}{(r_{13}+r_{14}+r_{34})^2}
            -\frac{2\big[x(r_{13}+r_{14}+r_{34})-(1-w)r_{13}\big]^2}{r_{34}^2(r_{13}+r_{14}+r_{34})^2(1-w)^2}\\
            &+\frac{1}{r_{34}(r_{13}+r_{14}+r_{34})}\left(\frac{1+x^2+(x+w)^2}{1-w}-\epsilon(1-w)\right).
    \end{split}
\end{equation}

\begin{mdframed}[backgroundcolor=gray!10, linewidth=0pt]
\begin{lstlisting}
In[1]:= B40 = BuildAntenna[B,4,0];

        tripleCollinearRules134 = {s13 -> lambda r13, s14 -> lambda r14,
            s34 -> lambda r34, s12 -> w sP2, s23 -> x sP2,
            s24 -> (1-w-x) sP2, s134 -> lambda (r13+r14+r34),
            s234 -> (1-w) sP2 + lambda r34,
            q2 -> sP2 + lambda (r13+r14+r34)};
        SeriesCoefficient[Together[B40 /. tripleCollinearRules134],
            {lambda, 0, -2}] // FullSimplify
\end{lstlisting}
\texttt{Out[1]:=}
    \begin{equation*}
        \frac{\begin{aligned}
            &-2r_{13}^2(w-1)^2+r_{13}r_{34}(w-1)\left(1+\epsilon(w-1)^2-w(w+2)\right)\\
            &+r_{14}r_{34}(w-1)\left(\epsilon(w-1)^2-1-w^2\right)+r_{34}^2w\left(1+\epsilon(w-1)^2-w^2\right)\\
            &-2(r_{13}+r_{14}+r_{34})\left[(w-1)(2r_{13}+r_{34}w)x+(r_{13}+r_{14}+r_{34}w)x^2\right]
        \end{aligned}}{r_{34}^2(r_{13}+r_{14}+r_{34})^2(w-1)^2}
    \end{equation*}
\end{mdframed}
Since $A_2^0=1$, this is the splitting function $P_\mathrm{NonIdent}^{(134)}$ of the equation above, written over the common denominator
$r_{34}^2(r_{13}+r_{14}+r_{34})^2(1-w)^2$.


\subsection{$C_4^0$}\label{subsec:localValidationC40}

The $C_4^0$ antenna is the identical-flavour four-quark antenna of the process $\gamma^*\to q\bar qq\bar q$,
built as the symmetrised direct--exchange interference. Unlike $B_4^0$, it has a single non-vanishing
limit, the triple-collinear limit $2_{\bar q}||3_q||4_{\bar q}$; the direct $3||4$ single-collinear region is
regular, as the check below confirms.

\subsubsection{Collinear $3||4$ Limit}

\begin{mdframed}[backgroundcolor=gray!10, linewidth=0pt]
\begin{lstlisting}
In[1]:= C40 = BuildAntenna[C,4,0];

        col34Rules = {s34 -> lambda, s13 -> z s1P, s14 -> (1-z) s1P,
            s23 -> z sP2, s24 -> (1-z) sP2, s134 -> s1P, s234 -> sP2,
            q2 -> s12 + s1P + sP2, Epsilon -> 0};

        FullSimplify[SeriesCoefficient[Together[C40 /. col34Rules],
            {lambda, 0, -1}]]
\end{lstlisting}
\texttt{Out[1]:=}
    \begin{equation*}
        0
    \end{equation*}
\end{mdframed}

\subsubsection{Triple-Collinear Limit}

For the $2_{\bar q}||3_q||4_{\bar q}$ limit, partons $2$, $3$, and $4$ become collinear together, leaving
parton $1$ as the sole hard remaining leg. As for $B_4^0$, we introduce the rescaled invariants
$r_{23}$, $r_{24}$, $r_{34}$, of order unity as $\lambda\to0$, via
\begin{equation}
    s_{23}\to\lambda r_{23},\quad s_{24}\to\lambda r_{24},\quad s_{34}\to\lambda r_{34},
\end{equation}
and the momentum fractions $w$, $x$, and $y\equiv1-w-x$ carried by partons $2$, $3$, and $4$
respectively within the recombined cluster, via
\begin{equation}
    s_{12}\to w\,s_{1P},\quad s_{13}\to x\,s_{1P},\quad s_{14}\to y\,s_{1P}.
\end{equation}
We also define $S_{234}\equiv r_{23}+r_{24}+r_{34}$.

At leading power, the four-parton antenna factorises onto the trivial LO antenna, $A_2^0\equiv1$, with an
explicit identical-particle factor of $\frac12$, as,
\begin{equation}
    \lim_{\lambda\to0}\lambda^2 C_4^0(1_q,3_q,4_{\bar q},2_{\bar q}) = \frac12 P_\mathrm{Ident}^{(234)}(w,x;r_{23},r_{24},r_{34})\,A_2^0(1,2),
\end{equation}
where the identical-flavour triple-collinear splitting function is a sum of two mirror-symmetric pieces \cite{Gehrmann-DeRidder:2005btv},
\begin{equation}
    \begin{split}
        P_\mathrm{Ident}^{(234)}(w,x;r_{23},r_{24},r_{34}) =\;&
            -\frac{2r_{23}/r_{34}+2}{S_{234}^2}-\frac{x(1+x^2)}{2r_{23}r_{34}(1-y)(1-w)}+\frac1{r_{34}S_{234}}\left(\frac{1+x^2}{1-y}+\frac{2x}{1-w}\right)\\
            &-\frac{2r_{34}/r_{23}+2}{S_{234}^2}-\frac{x(1+x^2)}{2r_{34}r_{23}(1-w)(1-y)}+\frac1{r_{23}S_{234}}\left(\frac{1+x^2}{1-w}+\frac{2x}{1-y}\right).
    \end{split}
\end{equation}

\begin{mdframed}[backgroundcolor=gray!10, linewidth=0pt]
\begin{lstlisting}
In[1]:= C40 = BuildAntenna[C,4,0];

        tripleCollinearRules234 = {s23 -> lambda r23, s24 -> lambda r24,
            s34 -> lambda r34, s12 -> w s1P, s13 -> x s1P,
            s14 -> (1-w-x) s1P, s234 -> lambda (r23+r24+r34),
            s134 -> (1-w) s1P + lambda r34,
            q2 -> s1P + lambda (r23+r24+r34), Epsilon -> 0,
            s123 -> (w+x) s1P + lambda r23,
            s124 -> (1-x) s1P + lambda r24};
        SeriesCoefficient[Together[C40 /. tripleCollinearRules234],
            {lambda, 0, -2}] // FullSimplify
\end{lstlisting}
\texttt{Out[1]:=}
    \begin{equation*}
        \frac{\begin{aligned}
            &-r_{23}^2\left(2w-(x-1)^2\right)(w+x-1)\\
            &+r_{23}\left[r_{24}(x-1)\left(1+w(x-1)+x(2x-1)\right)+r_{34}\left((x-1)^3-4w^2-4w(x-1)\right)\right]\\
            &+(r_{24}x-r_{34}w)\left(r_{24}-r_{34}+2r_{34}w+(r_{24}+r_{34})x^2\right)
        \end{aligned}}{2r_{23}r_{34}(w-1)(w+x)(r_{23}+r_{24}+r_{34})^2}
    \end{equation*}
\end{mdframed}
Since $A_2^0=1$, this is $\tfrac12P_\mathrm{Ident}^{(234)}$ of the equation above, written over a common denominator.


\subsection{Massive $A_3^0$}\label{subsec:localValidationMassiveA30}

For this massive antenna, described as,
\begin{equation}
    A_3^0\Big(1_Q,3_g,2_{\bar Q}\Big)
\end{equation}
where $Q$ denotes a quark of mass $m_q$, there is one infrared singular limit, for the soft gluon, 
and two mass-regulated quasi-collinear limits, for $Q||g$ and $\bar Q||g$. These two collinear limits 
are symmetric as $1\leftrightarrow2$ and thus, as before, the validation of one validates the other.
Relative to the massless case, it is relevant to note that there is no ordinary collinear pole
at fixed $m_q$.

\subsubsection{Soft-Gluon Limit}

This limit is computed by taking the gluon's momentum to zero, as $p_3\to0$. The antenna reduces to a product of 
an $A_2^0$ antenna and the massive eikonal factor as \cite{GehrmannDeRidder:2009fz},
\begin{equation}
    A_3^0(1_Q,3_g,2_{\bar Q})\xrightarrow{3_g\to0}\frac14S_{132}(m_q,m_q),\quad S_{132}(m_q,m_q) = \frac{2s_{12}}{s_{13}s_{23}}-\frac{2m_q^2}{s_{13}^2}-\frac{2m_q^2}{s_{23}^2}.
\end{equation}

\begin{mdframed}[backgroundcolor=gray!10, linewidth=0pt]
\begin{lstlisting}
In[1]:= mA30 = BuildAntenna[A,3,0,quarkMass->mQ];

        softRules = {s13 -> lambda s13, s23 -> lambda s23,
            s123 -> s12 + lambda (s13+s23),
            q2 -> s12 + 2 mQ^2 + lambda (s13+s23), Epsilon -> 0};
        SeriesCoefficient[Together[mA30 /. softRules],
            {lambda, 0, -2}] // FullSimplify
\end{lstlisting}
\texttt{Out[1]:=}
    \begin{equation*}
        \frac{s_{12}s_{13}s_{23}-m_q^2\left(s_{13}^2+s_{23}^2\right)}{2s_{13}^2s_{23}^2}
    \end{equation*}
\end{mdframed}
This is $\tfrac14\mathcal S_{132}(m_q,m_q)$ of the equation above, written over a common denominator.


\subsubsection{Collinear Limit and the Altarelli--Parisi Splitting Function}

To compute this limit, we start by replacing,
\begin{equation}
    \begin{gathered}
        p_1||p_3,\quad s_{13}\to 0,\quad m_q^2\to 0,\quad \frac{m_q^2}{s_{13}}\to \mu_{Qg}^2.
    \end{gathered}
\end{equation}
Here $\mu_{Qg}$ is some fixed quantity.
As usual, with $z$ representing the quark momentum fraction,
\begin{equation}
    p_1\to zp_{13},\quad p_3\to(1-z)p_{13}.
\end{equation}
$\epsilon$ is also set to zero.

The antenna factorised as \cite{GehrmannDeRidder:2009fz},
\begin{equation}
    A_3^0(1_Q,3_g,2_{\bar Q})\xrightarrow{1_Q||3_g}\frac1{4s_{13}}P_{Qg\to Q}(z,\mu_{Qg}^2),\quad
    \begin{gathered}
        P_{Qg\to Q}(z,\mu_{Qg}^2) = \frac{1+z^2}{1-z}-2\mu_{Qg}^2\\
        = \frac{1+z^2}{1-z},\quad\mu_{Qg}^2=0.
    \end{gathered}
\end{equation}

\begin{mdframed}[backgroundcolor=gray!10, linewidth=0pt]
\begin{lstlisting}
In[1]:= mA30 = BuildAntenna[A,3,0,quarkMass->mQ];

        quasi13Rules = {s13 -> lambda r13, s12 -> z sP2,
            s23 -> (1-z) sP2, s123 -> sP2 + lambda r13,
            q2 -> sP2 + lambda r13 (1+2 muQg2),
            mQ^2 -> lambda muQg2 r13, mQ^4 -> lambda^2 muQg2^2 r13^2,
            Epsilon -> 0};
        r13 SeriesCoefficient[Together[mA30 /. quasi13Rules],
            {lambda, 0, -1}] // FullSimplify
\end{lstlisting}
\texttt{Out[1]:=}
    \begin{equation*}
        \frac{1+z^2-2\mu_{Qg}^2(1-z)}{4(1-z)}
    \end{equation*}
\end{mdframed}
This is $\tfrac14P_{Qg\to Q}(z,\mu_{Qg}^2)$ of the equation above.

